% Standalone section; the parent report supplies amsmath/amsthm.
\section{A finite certificate for changing the decoder tree}
\label{sec:leaf-cfg-transfer}

Tree reassociation can shorten a decoder word or alter its zero runs,
but it does not automatically preserve this register language's behavior.
Break and continue count selected bits in the consumed path.  Their
numerical operands refer to the enclosing tree structure, so an ordinary
associativity claim for sequential programs is insufficient.
The finite leaf graph gives a direct semantic certificate.

Let $P$ be a finite full binary tree, with leaf-address set $L_P$ and the
literal source operations $\mathrm{CI}(r)$, $\mathrm{CD}(r)$,
$\mathrm{CB}(k)$, $\mathrm{CC}(k)$ and $\mathrm{CH}$.
Write $d_P(w)$ for the unique leaf reached by reading the finite word
$w$ followed by infinitely many zeros.  Write $u(w)$ for fixed-width
binary increment modulo $2^{|w|}$, including $u(\epsilon)=\epsilon$.
Let $s_b^k(w)$ remove the shortest suffix containing $k$ occurrences of
bit $b$, or remove the whole word if fewer than $k$ occur.  For $k=0$
it removes nothing.  These are precisely the source's
\texttt{incPath} and \texttt{stripPath} operations, before any
trailing-zero abbreviation.

Define the leaf control graph by the following transitions.  Every
edge carries the indicated register effect as well as its next leaf.
\[
\begin{array}{c|c|c}
\text{leaf operation at }\alpha&\text{register effect}&\text{next leaf}\\\hline
\mathrm{CI}(r)&R_r\gets R_r+1&d_P(u(\alpha))\\
\mathrm{CD}(r),\ R_r>0&R_r\gets R_r-1&d_P(u(\alpha))\\
\mathrm{CD}(r),\ R_r=0&\text{identity}&d_P(u(s_0^1(\alpha)))\\
\mathrm{CB}(k)&\text{identity}&d_P(u(s_0^k(\alpha)))\\
\mathrm{CC}(k)&\text{identity}&d_P(s_1^k(\alpha))\\
\mathrm{CH}&\text{halt}&\text{none}.
\end{array}
\]

\begin{lemma}[The selected leaf determines a source step]
If source paths $w$ and $w'$ select the same leaf $\alpha$ in $P$,
then their source executions from a common register file have the same
register effect, halting status, and next selected leaf.  This assertion
holds even if either path ends before the leaf or contains a suffix
beyond the leaf.
\end{lemma}
\begin{proof}
The decoder reads precisely the prefix $\alpha$ of length $|\alpha|$
from the zero-padded path.  The source definition applies
\texttt{consume} before executing the leaf.  Thus both inputs supply
the literal word $\alpha$ to \texttt{execute}.  Its six cases are exactly
the rows displayed above.  The tree is finite and full, so decoder fuel
exhaustion is impossible when its fuel bound is at least the tree's
maximum depth.  This is an essential distinction from treating every
\texttt{none} result in a partial decoder as a genuine halt.
\end{proof}

\begin{theorem}[Finite compiler-map certificate]
Let $P,Q$ be such trees and let $f:L_P\to L_Q$ be a total function.
Suppose the following finite conditions hold:
\begin{enumerate}
\item $f(d_P(\epsilon))=d_Q(\epsilon)$.
\item Increment and decrement leaves map to the same operation on the
same logical register.  Halt leaves map to halt leaves and only halt
leaves map to halt leaves.  Break and continue leaves may map to either
kind of unconditional identity-effect control leaf.
\item For every source leaf and each applicable guard, the next leaf
commutes with $f$.  In particular, both the zero and positive decrement
guards are checked separately.
\end{enumerate}
Then, for every initial register file $R$, the empty-path source runs
of $P$ and $Q$ have identical register files at every source step,
corresponding selected leaves under $f$, and the same halting status.
In particular their halting statements are equivalent in ZFC.
Neither injectivity nor surjectivity of $f$ is required.
\end{theorem}
\begin{proof}
Relate $(w,R)$ and $(v,S)$ when $R=S$ and
$d_Q(v)=f(d_P(w))$.  The first condition relates the initial states.
The selected-leaf lemma and the second condition give matching effects
and matching halt tags.  For a decrement, the common register file
selects the same guard.  The third condition then relates the two
successors.  Induction on the number of source steps gives the asserted
correspondence up to any halt.  Halt-tag reflection and determinism
give both directions of finite halting.  Extra target leaves cannot
cause an extra reachable behavior: initiality and successor closure
keep the target run in the image of $f$.
\end{proof}

The accompanying \texttt{leaf\_cfg.py} constructs this finite graph from
a complete site list.  It rejects duplicate addresses, non-prefix-free
addresses, incomplete binary trees, invalid opcodes, invalid operands,
incomplete map domains, nonexistent images, mismatched effects or halt
tags, and noncommuting successors.  Its default map pairs leaves in
left-to-right order, which is useful for reassociation retaining leaf
order.  An explicit map may also be supplied.

A fresh self-comparison of the inherited site list contains $6,460$
leaves and $8,391$ guarded cases.  It checks the implementation against
the complete finite input rather than a selection of instructions.
A three-leaf reassociation with two increments and a halt passes;
replacing its second increment with a control leaf is rejected.
These runs are checker validation, not a new optimized compiler result.
Binding a changed tree to its generated source and binary table remains
the responsibility of the compiler audit.

\subsection{Composing the compiler certificate with a changed binary machine}

For completeness, the final logical step is stated in a form that also
permits a new tape encoding or detector.  The positive-segment
composition principle is already present in the prior
\emph{Beyond 246} packet; the finite-time reflection proof below
spells out why unbounded macro duration creates no problem.

Let $\Gamma$ be the deterministic register graph and $M$ a deterministic
binary machine, with halting configurations made absorbing.  Let
$\mathcal R(p,c)$ relate an active graph configuration to a complete
binary configuration.  Assume a verified finite nonhalting blank prefix
ends at $c_0$ related to an admitted logical root $p_0$.  Assume every
active graph step has a binary simulation segment of length $k\geq1$
with the prescribed complete endpoint and no halt.  Assume each graph
halt has a finite binary halting segment.  All assumptions must be
proved on the chosen invariant, for arbitrary represented registers,
gaps and spare suffixes, not only on sampled values.

If the graph halts, induction on its finite execution concatenates
these segments and obtains a binary halt.  Conversely, use strong
induction on a finite binary halting time $t$ from a related
configuration.  If the graph is already at a halt leaf, it halts.  If
it takes an active step, its segment ends active after $k\geq1$ moves.
Since binary halting is absorbing, a halting time $t$ must satisfy
$t>k$.  Apply the induction hypothesis to the related endpoint with
halting time $t-k<t$.  This proves a finite graph halt.  No inference
from a long finite sample and no infinite choice of witnesses is
needed.  The finite blank prefix can be prefixed or removed using
determinism and its nonhalting certificate.

Finally, compose with the inherited equivalences
\[
 \operatorname{Halt}(\Gamma,p_0)
 \longleftrightarrow C_{254}(1)
 \longleftrightarrow\operatorname{Ref}(\mathrm{NativeAx})
 \longleftrightarrow\operatorname{Ref}(\mathrm{ZF})
 \longleftrightarrow\operatorname{Ref}(\mathrm{ZFC}).
\]
Classical negation gives the requested ZFC theorem.  A successful
finite compiler-map certificate does not replace the universal binary
segment proof or the new blank initialization.  Conversely, the latter
proofs need not preserve old decoder paths or old macro exit labels,
provided the complete related frames and positive progress are restored.
